% ===================== CAPÍTULO 5 - CONCLUSÕES ================================
\chapter{Conclusões}\label{chap:conclusoes}

\section{Conclusão}

O estágio permitiu participar no desenvolvimento e evolução de uma plataforma de seguimento marítimo utilizada em contexto empresarial real, envolvendo a integração de fornecedores externos, o desenvolvimento de funcionalidades de frontend e backend e a adaptação de componentes críticos da arquitetura da solução.

Os objetivos definidos para o estágio foram, de forma geral, cumpridos. Destacam-se a integração do fornecedor SeaRates, a migração do sistema de autenticação para Keycloak, a correção de problemas relacionados com a representação geográfica dos trajetos marítimos e a evolução do motor de estatísticas da plataforma.

Para além dos resultados técnicos alcançados, a experiência proporcionou contacto direto com desafios característicos do desenvolvimento de software em ambiente empresarial, nomeadamente a necessidade de compreender sistemas existentes, adaptar prioridades em função das necessidades do produto e tomar decisões com impacto na manutenibilidade e evolução futura da solução.

Em termos pessoais e profissionais, o estágio contribuiu para consolidar competências de desenvolvimento de software, autonomia, comunicação técnica e trabalho em equipa, constituindo uma etapa relevante na transição entre o contexto académico e a prática profissional.

\section{Reflexão Crítica}

\subsection{Atividades Desenvolvidas e Competências Adquiridas}

Ao longo do estágio, foram adquiridas e consolidadas competências técnicas e pessoais diretamente relacionadas com as atividades realizadas. O trabalho na \textit{Tracking \gls{api}} proporcionou o aprofundamento da componente de desenvolvimento \textit{backend} com .NET - serviços \textit{REST}, integração com \glspl{api} externas e mapeamento de \textit{\glspl{dto}} - complementada pelo uso do \gls{efcore} no acesso a dados em SQL Server. O desenvolvimento de \textit{frontend} com React, abrangendo a criação de interfaces, a integração com \glspl{api} e a gestão de estado, completou o leque de competências técnicas desenvolvidas.

Foram ainda utilizadas ferramentas de desenvolvimento profissional como Git, Azure DevOps~\cite{azuredevops}, Swagger~\cite{swagger} e as ferramentas de desenvolvimento do navegador, para além da adaptação a arquiteturas de \textit{software} existentes e consolidadas.

O regime de trabalho parcial, com disponibilidade condicionada pelo horário letivo, exigiu uma evolução rápida na autonomia de resolução de problemas, na aprendizagem contínua e na gestão do tempo - capacidades que o contexto do estágio naturalmente estimulou. A comunicação em contexto de equipa, incluindo a discussão de soluções técnicas e o esclarecimento de dúvidas, bem como a capacidade de integrar e compreender código existente, foram também aspetos relevantes do crescimento profissional alcançado.

\subsection{Estratégias de Trabalho Adotadas}

A execução do projeto assentou numa metodologia ágil adaptada à realidade corporativa da Devlop. O fluxo de trabalho foi estruturado em iterações semanais, com definição prévia de objetivos operacionais, em moldes semelhantes aos \textit{sprints} de desenvolvimento. No final de cada ciclo, realizavam-se sessões de alinhamento com o orientador da entidade de acolhimento para a validação das decisões de arquitetura e a demonstração dos resultados alcançados.

A estratégia de integração foi também acelerada pela partilha diária de conhecimento com o colega de equipa, Sandro Cunha, o que facilitou a assimilação da base de código pré-existente e a mitigação dos constrangimentos associados à curva de aprendizagem das ferramentas.

\subsection{Tecnologias e Ferramentas Utilizadas}

A \autoref{tab:tecnologias} (\autoref{sec:tecnologias}) sistematiza o conjunto de tecnologias e ferramentas utilizadas ao longo do estágio, organizadas por categoria e por âmbito de aplicação. A escolha da \textit{stack} tecnológica foi determinada pelo contexto da empresa, que privilegia o ecossistema Microsoft (.NET, SQL Server) no \textit{backend} e React no \textit{frontend}, complementada por ferramentas de integração contínua (Azure DevOps), autenticação (Keycloak) e produtividade (Secção~\ref{sec:ia}). Para teste e validação, recorreu-se a um ecossistema consolidado de \textit{mocking} (Prism) e de testes automatizados (xUnit, Moq, FluentAssertions), assegurando a qualidade e a fiabilidade da solução desenvolvida.

\subsection{Planeamento Previsto e Executado}

A execução do projeto exigiu adaptações face ao plano de trabalho inicialmente delineado no \autoref{chap:planeamento}. Embora o objetivo não passasse por implementar um sistema de \textit{tracking} de raiz, a complexidade de adicionar um novo fornecedor ao ecossistema existente exigiu um esforço de análise e desenvolvimento substancial, distribuído ao longo de 16 semanas. O \anexoref{anx:execucao-real} (\autoref{tab:execucao-real}) apresenta o registo detalhado da execução real do estágio, cruzando as atividades efetivamente desenvolvidas com a respetiva calendarização.

Comparativamente ao \gls{pit}, verifica-se que as Fases~1 e~2 decorreram conforme o previsto, totalizando 56 horas nas duas primeiras semanas. A Fase~3, embora tenha mantido a sua duração global, foi concluída antecipadamente (72 horas em três semanas, face às 160 horas e quatro semanas planeadas), beneficiando da utilização de ferramentas de simulação que agilizaram o desenvolvimento e os testes. A Fase~4, relativa ao \textit{frontend}, registou o maior desvio: as 320 horas inicialmente estimadas foram ajustadas para 168 horas reais, distribuídas ao longo de seis semanas, registou-se uma redução do número de horas efetivamente necessárias face ao inicialmente previsto, acompanhada por uma delimitação mais precisa do âmbito funcional (\autoref{tab:execucao-real}).

A Fase~5, originalmente dedicada à integração com o i-Forwarder, foi reorientada para incluir também a preparação do ambiente \gls{4gl} (\anexoref{anx:4gl}) e a resolução de problemas de paginação, totalizando 68 horas em três semanas. Por fim, a Fase~6 alargou o seu âmbito para acomodar a migração para Keycloak e a refatoração do motor de estatísticas, ocupando 52 horas em duas semanas, conforme sistematizado na \autoref{tab:comparacao-horas}.

A \autoref{tab:comparacao-horas} (\autoref{sec:cumprimento}) apresenta a comparação detalhada entre as horas planeadas e as executadas, evidenciando uma redução global de 38\,\% face ao \gls{pit} inicial.

\subsection{Dificuldades Encontradas}

Ao longo do estágio, foram identificadas dificuldades de natureza organizacional e técnica. No plano organizacional, a gestão do tempo constituiu um desafio relevante, uma vez que o estágio decorreu em regime parcial durante a maior parte do semestre. A disponibilidade esteve condicionada pelo horário letivo da licenciatura, que ocupava as segundas e terças-feiras quase na totalidade (\anexoref{anx:horario}, \autoref{fig:horario}), exigindo uma rigorosa organização de tarefas e um foco redobrado nas atividades críticas nos restantes dias úteis.

No plano técnico, uma das principais dificuldades esteve relacionada com o desenvolvimento de \textit{frontend}, área na qual havia menor experiência prévia. A assimilação do modelo React disponibilizado pela equipa apresentou-se como um desafio, dada a sua complexidade e arquitetura específica.

Como forma de ganhar tração técnica, adotou-se uma abordagem exploratória: dedicaram-se três dias úteis à implementação isolada de uma interface \textit{\gls{crud}} completa com React Admin~\cite{reactadmin}, o que foi fundamental para interiorizar a gestão de estado e a integração com a \gls{api}. Com o tempo e o apoio pontual dos colegas, esta barreira inicial foi superada.

A complexidade da arquitetura de \textit{backend} e a necessidade de integrar um novo fornecedor num sistema já consolidado (\anexoref{anx:arquitetura}, \autoref{fig:arquitetura}) representaram outro desafio, ultrapassado através da análise detalhada do código e de testes progressivos. Na reta final, a migração para Keycloak e a resolução da anomalia do antimeridiano exigiram um esforço analítico profundo, mas revelaram-se duas das experiências de aprendizagem mais marcantes do estágio.

A implementação do motor de estatísticas também apresentou um desafio arquitetural inesperado. A lógica inicial, baseada em agregações pré-calculadas, apresentava fragilidades e, após uma primeira tentativa de refatoração, o sistema tornou-se excessivamente complexo. A principal dificuldade - e consequente aprendizagem - foi reconhecer esse excesso e redesenhar a solução para uma abordagem mais simples, consultando diretamente os dados da \gls{bd}.

Esta simplificação, orientada pela manutenibilidade, trouxe um benefício adicional: ao operar com dados em bruto, o sistema ganhou maior granularidade, facilitando a implementação da funcionalidade de exportação para o Excel.

\section{Trabalho Futuro}

Como propostas para desenvolvimentos futuros e para a otimização do ecossistema desenvolvido, identificam-se três linhas de ação prioritárias.

Em primeiro lugar, recomenda-se a implementação de um mecanismo de armazenamento em \textit{cache} e de persistência local para os dados da SeaRates. Dado que esta \gls{api} opera num modelo sem estado, a criação de uma \textit{cache} na \gls{bd} interna da Devlop permitiria reter o histórico de estados e trajetórias dos navios por períodos mais longos, reduzindo o volume de consultas redundantes à \gls{api} externa e minimizando os custos operacionais associados.

Em segundo lugar, sugere-se retomar a integração com o sistema i-Forwarder. Uma vez que o ambiente de desenvolvimento em Four Js Genero foi preparado e ocorreram os alinhamentos iniciais na fase final do estágio, a criação de \textit{triggers} automáticos entre o ecossistema \gls{4gl} do i-Forwarder e a \textit{Tracking \gls{api}} permitiria que qualquer atualização de estado se refletisse instantaneamente em alertas operacionais nos ecrãs dos consultores, eliminando a necessidade de sincronização manual.

Por último, propõe-se uma reavaliação do motor de estatísticas. A decisão de remover as agregações pré-calculadas em favor da consulta direta à \gls{bd} simplificou o código atual, mas, à medida que a plataforma escalar e o volume de histórico aumentar, a execução destas consultas em tempo real poderá gerar estrangulamentos de desempenho.

Para garantir a escalabilidade sem reintroduzir lógicas frágeis, sugere-se a implementação de soluções ao nível da \gls{bd} - como as vistas materializadas - ou a introdução de uma camada de \textit{cache} dedicada às métricas que não exijam atualização imediata.

